\subsection{NN-G034: the exact optimum with one fixed first-cell realization is 2744}
\paragraph{Statement and class.}
In the symmetric closed-table presentation, impose only
\[
 \delta_1=D(D-3),\qquad G_1=2D^2(D-2),\qquad U_1=D.   \tag{G034.1}
\]
Allow all remaining actual coefficient polynomials to vary,
with arbitrary degrees and contents, and allow rational
auxiliary $\Lambda_i$ and nonconstant ratios.
Among effective associative lifts in this class, the exact
minimum is $2744$ cancelling pairs. There is exactly one
minimizing set of actual base polynomials, namely NN-G033.
This is an exact class theorem, not the unrestricted
$1000$-position optimum and not the optimum over all closed
tables with their first-cell data free.

\paragraph{Forced parameters and the first-cell optimum.}
Put $C=D-1$ and $H=2D(D-2)$, so $G_1=DH$.
The identities $U_1G_2=W(G_1-D\delta_1)$ and
$WS=G_1G_2$ force
\[
                  G_2=DWC,\qquad S=D^2CH.            \tag{G034.2}
\]
Thus $(g_1,g_2,s)=(64,12w,768)$ at $D=4$.
The first $V$ identity forces
\[
 A_1=D^2H-DCV_1,\qquad a_1=256-12v_1.
\]
Since $v_1$ is even and $a_1\ge8$, one has $v_1\le20$.
The first-cell contribution satisfies
\[
 20v_1+14a_1=3584-148v_1\ge624,
\]
with equality exactly when $v_1=20,a_1=16$.

There is exactly one allowable $V_1$ at this equality.
It dominates the trace character $5D-4=(6,5)$,
has the same parity, and has surplus mass four.
Since $A_1$ is effective and its first term $D^2H$
has degree four, $V_1$ has degree at most two.
The only possibilities are
\[
 5D,\qquad7D-8,\qquad5D-4+2(q^2+q^{-2})=2D^2-3D.
\]
The respective central coefficients of $A_1$ are
$-10,-6,2$. The last choice gives
$A_1=D^2(D-3)=(2,3,3,1)$ and satisfies all conditions;
the first two fail effectivity.

\paragraph{A complete mass reduction.}
For a competitor with score at most $2744$, the nonunit
mass is $T\le9912$. Formula (G031.1) and the preceding
first-cell bound give
\[
 \begin{aligned}
 T&\ge2856+440w+20u_2+20v_2+26d_2+14a_2,\\
 u_2&=\frac{w(3w-d_2)}{16},\qquad
 v_2=\frac{w(48w-a_2)}{64}.                         \end{aligned} \tag{G034.3}
\]
The omitted amount is exactly $148(20-v_1)\ge0$.
Using $u_2\ge4,v_2\ge16,d_2\ge4,a_2\ge16$ gives
$T\ge3584+440w$, hence $w\le14$.
NN-G031 rules out $w=4,6$ at this score.
Eliminating $u_2,v_2$ in (G034.3) gives
\[
 T\ge2856+440w+\frac{75}{4}w^2
       +\left(26-\frac{5w}{4}\right)d_2
       +\left(14-\frac{5w}{16}\right)a_2.           \tag{G034.4}
\]
For $w=12,14$ both final coefficients are positive,
and even the first three terms are at least
$10836>9912$. Thus only $w=8$ or $w=10$ remain.

Write $T_j=q^j+q^{-j}$, an integral polynomial in $D$.
Every effective symmetric $W\equiv D\pmod2$
dominating $D$ has the following exhaustive small-mass forms:
\[
 \begin{array}{c|l}
 w=8&D+4\quad\hbox{or}\quad D+2T_j\quad(j\ge1),\\
 w=10&D+6\quad\hbox{or}\quad D+2+2T_j\quad(j\ge1).
 \end{array}                                                       \tag{G034.5}
\]
Indeed the even surplus has mass four or six; a noncentral
surplus pair costs four, leaving only the displayed central
remainder. Each polynomial in (G034.5) is primitive in
$\mathbb Z[D]$, since its Laurent coefficients have gcd one.

\paragraph{An equal-mass divisibility lemma.}
Suppose $W,U$ both belong to the same row of (G034.5),
and $U=WR$ for $R\in\mathbb Z[D]$ with $R\equiv1\pmod2$.
Then $U=W$.
If $\deg W\ge2$, its leading coefficient is two.
A nonconstant quotient would require $\deg U\ge2$,
also with leading coefficient two; thus $R$ would have
leading coefficient one, contradicting $R\equiv1$
in positive degree. A constant quotient has $R(4)=1$
and is one. For the degree-one choice with leading
coefficient three, a nonconstant integral quotient would
have to have leading coefficient $2/3$, impossible.

It remains to consider $W=D+\ell$ with $\ell=4$ or six.
A nonconstant quotient requires
$U=D+\ell-4+2T_j$ with $j\ge2$.
At $D=-\ell$, divisibility would give $T_j(-\ell)=2$.
This is impossible. To see it directly, set $x=\ell+2\ge6$;
the positive sequence
$a_j=(-1)^jT_j(-\ell)$ has $a_0=2,a_1=x$ and
$a_{j+1}=xa_j-a_{j-1}$. Induction gives
$a_{j+1}>a_j>2$ for $j\ge1$.
This completes the lemma without bounding $j$.

\paragraph{The complete $w=8$ branch.}
Each $W$ in its row of (G034.5) is coprime over
$\mathbb Q[D]$ to $D(D-2)$: at $D=0$ its value is
$4$ or $\pm4$, and at $D=2$ its value is among
$6,2,-2$. The two second-cell identities therefore
force
\[
 U_2=WR,\qquad V_2=WZ\qquad(R,Z\in\mathbb Z[D]).
\]
Here polynomial divisibility follows first over
$\mathbb Q[D]$ and then integrally by primitivity.
Cancelling $D$ modulo two gives $R\equiv Z\equiv1$.
Formula (G034.3) shows $4\le u_2\le10$, so the
positive odd integer $R(4)=u_2/8$ must be one.
The equal-mass lemma gives $U_2=W$, and then
\[
                 \delta_2=WC-H,\qquad d_2=8.         \tag{G034.6}
\]
Likewise $16\le v_2\le46$, so $Z(4)$ is three or five.
The corresponding $a_2$ are $192$ and $64$.

For $W=D+4$, (G034.6) has negative leading coefficient,
so it is not effective. For $W=D+2T_j$, one has
\[
             \delta_2=-D(D-3)+2(D-1)T_j.
\]
Domination of $D(D-3)$ forces $j\le2$: when $j\ge3$
the product $(D-1)T_j$ has no coefficient at $q^1$,
whereas $D(D-3)$ does. Thus only $W=3D-4$ and
$W=D+2T_2$ remain.

The second $V$ identity is now
\[
              A_2=D^2\bigl(WC-2(D-2)Z\bigr).         \tag{G034.7}
\]
If $W=3D-4$, effectivity of $V_2=WZ$ makes the leading
coefficient of a nonconstant $Z$ positive. Its parity
makes that coefficient even. For $\deg Z=1$,
the leading coefficient inside (G034.7) is at most
$3-4<0$, and higher degrees also give a negative
leading coefficient. Thus $Z$ must be constant.
But then $V_2$ has degree one and cannot dominate
the trace coefficient two at $q^2$. This excludes
the degree-one choice.

For $W=D+2T_2=(2,1,2)$, the same leading-coefficient
argument gives $\deg Z\le1$: degree two would have
leading coefficient at least two, against the leading
coefficient two of $WC$.
Write $Z=(z-2a)+a(q+q^{-1})$, where $z=Z(4)\in\{3,5\}$.
Here $a\ge0$ is even, while the central coefficient
is allowed to be negative. The central coefficient
of $V_2=WZ$ is $2z-2a$ and must be at least six.
If $z=3$, this forces $a=0$ and $Z=3$.
If $z=5$, it permits only $a=0,2$.
The central coefficient of $A_2$ is
$70-16z+4a$, which is negative in both $z=5$ cases.
Consequently only $Z=3$ survives. Its $a_2=192$
gives, from (G034.4),
\[
                  T\ge7576+16\cdot8+\frac{23}{2}\cdot192
                    =9912.                         \tag{G034.8}
\]
Equality forces the unique first-cell optimum already
proved, and gives exactly the parameters of NN-G033.

\paragraph{The complete $w=10$ branch: scalar constraints.}
Here (G034.3) reads
\[
 u_2=\frac58(30-d_2),\qquad
 v_2=\frac5{32}(480-a_2).
\]
Evenness and positivity force $d_2\equiv14\pmod{16}$
and $d_2<30$, hence $d_2=14,u_2=10$.
They also force $a_2\equiv32\pmod{64}$.
The choice $a_2=32$ gives the lower bound $T\ge9668$,
whereas the next choice $a_2=96$ already gives
$T\ge10364>9912$. Thus a competitor must have
\[
                  d_2=14,\quad u_2=10,\quad
                  a_2=32,\quad v_2=70.               \tag{G034.9}
\]

\paragraph{Coprime polynomials in the $w=10$ branch.}
No $W$ in this row is divisible by $D$.
It is divisible by $D-2$ exactly when
$W=D+2+2T_j$ with $j\equiv2\pmod4$, since
$T_j(2)$ is zero for odd $j$ and
$2(-1)^{j/2}$ for even $j$.
First suppose $W$ is coprime to $D(D-2)$.
Then $U_2=WR,V_2=WZ$ integrally, with
$R\equiv Z\equiv1$, and $R(4)=1$.
The equal-mass lemma gives $U_2=W$ and
$\delta_2=WC-H$.

The choice $W=D+6$ again gives a negative leading
coefficient in $\delta_2$. For
$W=D+2+2T_j$, domination of $\delta_2$ requires
\[
                       (D-1)(1+T_j)\ge D(D-3).
\]
The coefficient at $q^2$ forces $j\le3$.
The coprime possibilities are $j=1$ and $j=3$.
For $j=1$, $W=3D-2$ and the same leading-coefficient
argument in (G034.7) forces $Z$ constant, so $V_2$
cannot dominate its required $q^2$ coefficient.

For $j=3$, effectivity in (G034.7) forces
$\deg Z\le2$: a degree-three leading coefficient
would be positive even and produce a negative leading
term, and larger degrees are also impossible.
Write $Z=(z_0,z_1,z_2)$ in symmetric Laurent notation.
By (G034.9), $Z(4)=7$.
The coefficient of $q^2$ in $D^2WC$ is $54$.
Since $G_1=(16,14,8,2)$, the coefficient of $q^2$
in $G_1Z$ is
\[
        8z_0+16z_1+16z_2=8Z(4)=56.
\]
Thus $A_2=D^2WC-G_1Z$ has coefficient $-2$ at $q^2$,
contradicting effectivity. This excludes every coprime
$w=10$ choice, including all unbounded locations $j$.

\paragraph{The $D-2$ factor in the $w=10$ branch.}
It remains to consider $W=D+2+2T_j$ with
$j\equiv2\pmod4$. Write $W=(D-2)B$.
The polynomial $B$ is primitive, $B(0)=-3$, and
$B(2)=1$: for even $j$, $T_j$ is an even polynomial
in $D-2$, so $W'(2)=1$.
Thus $B$ is coprime to $D(D-2)$; modulo two,
$B\equiv1$. The identities now force
\[
 U_2=BR,\quad V_2=BZ,\qquad
 R\equiv Z\equiv D\pmod2,
 \qquad R(4)=2,\quad Z(4)=14.                       \tag{G034.10}
\]
The leading coefficient of $B$ is two and
$\deg B=j-1$. If $R$ were constant, its parity
would contradict $R\equiv D$.
Otherwise $\deg U_2\ge j\ge2$, so the mass-ten
classification gives leading coefficient two for $U_2$.
Hence $R$ has leading coefficient one, and its parity
forces $\deg R=1$. The only mass-ten polynomial of
degree $j$ in (G034.5) is $W$ itself. Therefore
$U_2=W$, $R=D-2$ and $\delta_2=WC-H$.
The preceding $q^2$ domination argument gives $j\le3$;
thus necessarily $j=2$.

We have reached
\[
 W=2D^2-7D+6=(D-2)(2D-3),\qquad B=2D-3.
\]
The second $V$ identity reduces to
\[
                    A_2=D^2(WC-2Z).                 \tag{G034.11}
\]
Effectivity forces $\deg Z\le2$. Indeed its leading
coefficient is positive when nonconstant; degree three
would make it positive even by (G034.10), exceeding
the leading coefficient two of $WC$ after the factor
two in (G034.11), and higher degrees also fail.
Write $Z=(z_0,z_1,z_2)$. Its parity says that $z_2$
is even, $z_1$ is odd, and $z_0$ is even; its mass says
$z_0=14-2z_1-2z_2$.
The relevant coefficients of (G034.11) are
\[
 \begin{array}{c|l}
 q^1&-25+2z_1+8z_2,\\
 q^2&32-4z_1-8z_2,\\
 q^4&11-2z_2.
 \end{array}
\]
Domination of their trace targets $3,2,1$ implies
\[
 z_1+4z_2\ge14,\qquad z_1+2z_2\le7,\qquad z_2\le5.
\]
The first two inequalities force $z_2\ge7/2$.
Evenness then gives $z_2=4$, and oddness gives
$z_1=-1$, hence $z_0=8$.
But the central coefficient of $V_2=BZ$ is then
$z_0+4z_1=4$, below its required value six.
This excludes the final $w=10$ branch.

The $w=8$ equality realization of NN-G033 is therefore
the unique minimizer in (G034.1), proving the theorem.
\hfill$\square$

\paragraph{Geometric translation and remaining gate.}
Fixing (G034.1) fixes the first-cell coefficient object,
its bottom action character, and one gluing product in
the NN-G010 category of graded sign-representation kernels.
Among all effective associative completions in this closed
presentation, the least total excess of dimension over
absolute trace is exactly $2\cdot2744$, achieved by the
single set of coefficient characters of NN-G033.
The proof allows arbitrary grading locations in the other
objects; the mass-ten $D-2$ factor is treated by an actual
integral divisibility argument, not discarded as generic.

The finite-character evaluations $D=0$ and $D=2$
are restrictions at $q=-1$ and $q=i$, under SGA 3,
I, 4.7.2--4.7.3. They are not assertions about Frobenius
eigenvalues. The monoidal composition dictionary is
archive cote 2, scan 17, via NN-G010.
Any smaller closed-table score must change at least one
of the three first-cell data in (G034.1); other table
shapes remain outside this class. This is a strict
remaining gate, with no canonical-basis, rank-four,
PH0/PH1 or purity conclusion.

